\documentclass[11pt]{article}

\usepackage[T1]{fontenc}
\usepackage[utf8]{inputenc}
\usepackage{lmodern}
\usepackage{amsmath}
\usepackage{amssymb}
\usepackage{enumitem}

\title{Option Sets, Agreement, Conflict, Precedence, and Difference}
\author{}
\date{}

\begin{document}

\maketitle

\section{Definitions}

Define:

\begin{enumerate}[label=\arabic*.]
    \item a \textbf{type} \(\tau\);
    \item a \textbf{key} \(k\), a unique identifier with \(k : \tau\);
    \item a \textbf{value} \(v\): an optional datum of arbitrary type that an option may carry;
    \item an \textbf{option} \(o\): a key \(k(o)\) together with an optional value \(v(o)\). The key is mandatory; the value is not.
    \item an \textbf{option set} \(O = \{o_1, o_2, \ldots, o_n\}\): a finite set of options in which the key is unique --- i.e. the key map \(k(\cdot)\) is injective on \(O\) (\(o_i \neq o_j \Rightarrow k(o_i) \neq k(o_j)\)).
\end{enumerate}

The \textbf{key set} of an option set is

\[
K(O) = \{\, k(o) \mid o \in O \,\}.
\]

Because keys are unique, \(k(\cdot)\) is a bijection between \(O\) and \(K(O)\). We write \(O[k]\) for the unique option carrying key \(k\):

\[
O[k] := \text{the unique } o \in O \text{ with } k(o) = k,
\qquad k \in K(O).
\]

Two options are \textbf{identical} (\(o_1 = o_2\)) iff they have the same key and the same value:
\(k(o_1) = k(o_2)\) and \(v(o_1) = v(o_2)\). (If an option carries several data rather than a single value, identity is componentwise over all of them.)

\section{Comparing Option Sets}

Two option sets are \textbf{equal} iff they contain the same options:

\[
O_1 = O_2
\iff
K(O_1) = K(O_2)
\;\text{and}\;
\forall k \in K(O_1),\; O_1[k] = O_2[k].
\]

By extensionality this is just mutual containment; equal cardinality follows as a consequence rather than a separate requirement.

Two option sets are \textbf{in agreement} iff, for every key occurring in both, the corresponding options are identical:

\[
O_1 \sim O_2
\Longleftrightarrow
\forall k \in K(O_1) \cap K(O_2),\; O_1[k] = O_2[k].
\]

\textbf{Remark.} \(\sim\) is reflexive and symmetric but \emph{not} transitive: it is a tolerance (compatibility) relation, not an equivalence. Disjoint key sets agree vacuously, so \(O_1 \sim O_2\) and \(O_2 \sim O_3\) do not imply \(O_1 \sim O_3\). In particular, \(\sim\) does not partition option sets into equivalence classes and must not be used to group or deduplicate them.

Two option sets are \textbf{in conflict} iff there is at least one shared key whose options are not identical:

\[
O_1 \not\sim O_2
\Longleftrightarrow
\exists k \in K(O_1) \cap K(O_2)\; \text{such that}\; O_1[k] \neq O_2[k].
\]

Conflict is exactly the negation of agreement, so the two are complementary.

\textbf{Remark.} Conflict requires a \emph{shared} key. Options that differ but whose keys differ do not conflict: if \(K(O_1) \cap K(O_2) = \varnothing\) the sets are vacuously in agreement. For example, \(O_1 = \{(\mathsf{foo}, 42)\}\) and \(O_2 = \{(\mathsf{bar}, \mathsf{hello})\}\) are never in conflict.

\subsection*{Union and intersection}

The \textbf{intersection} keeps the options on which the two sets coincide:

\[
O_1 \cap O_2 := \{\, o \in O_1 \mid k(o) \in K(O_2) \;\wedge\; O_2[k(o)] = o \,\}.
\]

It is empty whenever the sets share no key (so for the example above, \(O_1 \cap O_2 = \varnothing\)).

The ordinary \textbf{union}

\[
O_1 \cup O_2 := \{\, o \mid o \in O_1 \;\vee\; o \in O_2 \,\}
\]

is a valid option set \emph{iff} \(O_1 \sim O_2\); otherwise a shared key would carry two distinct options, violating uniqueness. When defined, \(K(O_1 \cup O_2) = K(O_1) \cup K(O_2)\). On disjoint sets the union always exists and coincides with the precedence-resolved union \(\oplus\) of \S3 in either order. In short, \(\cup\) is the agreement-restricted total of \(\oplus\); \(\oplus\) extends it to conflicting sets by electing a winner.

\section{Precedence}

\[
O_2 \succ O_1
\]

means that \(O_2\) has precedence over \(O_1\).

Define a \textbf{precedence-resolved union}:

\[
O_1 \oplus O_2.
\]

This may be read as:

\begin{quote}
\(O_1\) is overridden by \(O_2\).
\end{quote}

Equivalently:

\begin{quote}
Combine \(O_1\) and \(O_2\), giving precedence to \(O_2\).
\end{quote}

Formally:

\[
K(O_1 \oplus O_2) = K(O_1) \cup K(O_2),
\]

and:

\[
(O_1 \oplus O_2)[k]
=
\begin{cases}
O_2[k], & k \in K(O_2), \\
O_1[k], & k \in K(O_1) \setminus K(O_2).
\end{cases}
\]

That says:

\begin{itemize}
    \item keep every option whose key appears in either set;
    \item if only \(O_1\) has the key, use \(O_1\)'s option;
    \item if \(O_2\) has the key, use \(O_2\)'s option;
    \item therefore, when both sets share a key, \(O_2\) overrides \(O_1\).
\end{itemize}

\textbf{Proposition (commutativity \(\iff\) agreement).}
\(O_1 \oplus O_2 = O_2 \oplus O_1\) if and only if \(O_1 \sim O_2\).

\emph{Proof sketch.} Off the overlap each side is forced to the unique present option, so order is irrelevant there. On \(K(O_1) \cap K(O_2)\), \(O_1 \oplus O_2\) yields \(O_2[k]\) while \(O_2 \oplus O_1\) yields \(O_1[k]\); these agree for every shared \(k\) iff \(O_1[k] = O_2[k]\) throughout, i.e. iff \(O_1 \sim O_2\). \(\square\)

\textbf{Remark.} \(\oplus\) is associative with identity \(\varnothing\), so \((\text{option sets}, \oplus, \varnothing)\) form a monoid --- non-commutative in general, by the proposition. A left fold \(O_{\mathrm{base}} \oplus O_1 \oplus \cdots \oplus O_n\) is therefore well-defined and order-consistent (``later wins'').

\section{Difference}

Fix a \textbf{value projection} \(\pi\) mapping an option to the carrier being compared,

\[
\pi(o) \in \{\, \bot \,\} \cup \{\text{values}\},
\]

where \(\bot\) denotes ``no value of the requested kind''. \(\pi\) selects \emph{which} aspect of an option is compared --- for instance its declared value, a default, or an effective ``value-or-default''. Carrier comparison is structural, with \(\bot \neq v\) for any value \(v\) and \(\bot = \bot\).

Over the shared keys, the \textbf{changed set} under \(\pi\) is

\[
\operatorname{changed}_{\pi}(O_a, O_b)
=
\{\, k \in K(O_a) \cap K(O_b) \mid \pi(O_a[k]) \neq \pi(O_b[k]) \,\},
\]

and the key-set delta is

\[
\operatorname{added}(O_a, O_b)   = K(O_b) \setminus K(O_a),
\qquad
\operatorname{removed}(O_a, O_b) = K(O_a) \setminus K(O_b).
\]

The \textbf{full difference} is their union:

\[
\operatorname{diff}_{\pi}(O_a, O_b)
=
\operatorname{added} \,\cup\, \operatorname{changed}_{\pi} \,\cup\, \operatorname{removed}.
\]

\textbf{Relation to \S2.} \(\operatorname{changed}_{\pi}\) alone characterizes agreement \emph{under} \(\pi\): it is empty iff the two sets agree on their overlap (with \(\pi = \mathrm{id}\) this is exactly \(O_a \sim O_b\)). The full \(\operatorname{diff}_{\pi}\) is empty iff the sets are equal under \(\pi\) --- same keys, same carriers everywhere --- so equality, unlike agreement, additionally requires the added and removed parts to vanish.

\textbf{Exemptions.} To ignore selected keys or options, filter the \emph{result} with a predicate \(\varepsilon\) rather than threading it through the comparison:

\[
\operatorname{diff}_{\pi,\varepsilon}(O_a, O_b)
=
\{\, k \in \operatorname{diff}_{\pi}(O_a, O_b) \mid \neg \varepsilon(k) \,\}
\]

(where \(\varepsilon\) may consult \(O_a[k]\) and \(O_b[k]\)). Keeping projection (\emph{what} is compared) and exemption (\emph{which} keys are reported) on separate axes leaves both independently composable; \(\varepsilon\) need not be part of the difference operator itself.

\end{document}
